% --- Archon physics formalization source begin ---
% archon:physics
% archon:covers PhyXMiniProblems/problem_phyx_mini_0862.lean
% archon:source-report reports/phyx_mini/problem_phyx_mini_0862.source.json

\chapter{Physics problem phyx\_mini\_0862}
\label{ch:PhyXMiniProblems_problem_phyx_mini_0862}

\paragraph{Problem source.}
\#\# Physical scenario

Two point charges \textbackslash{}( q\_a \textbackslash{}) and \textbackslash{}( q\_b \textbackslash{}) are located on the \textbackslash{}( x \textbackslash{})-axis at \textbackslash{}( x = a \textbackslash{}) and \textbackslash{}( x = b \textbackslash{}). figure is a graph of \textbackslash{}( E\_x \textbackslash{}), the \textbackslash{}( x \textbackslash{})-component of the electric field.

\#\# Question

What is the ratio \textbackslash{}( \textbackslash{}left| q\_a / q\_b \textbackslash{}right| \textbackslash{})?

\#\# Answer choices

- A: A: 1.2

- B: B: 1.4

- C: C: 2

- D: D: 1

\#\# Figure caption (auxiliary; use the image as primary evidence)

The image shows a mathematical graph with the following features:

- A continuous green curve that forms a semicircular shape opening downward between two vertical dashed lines.
- The two vertical dashed lines are labeled with "a" and "b" at their intersection with the x-axis, indicating specific points or boundaries.
- The x-axis is labeled with "x."
- The curve and dashed lines create two closed regions that are symmetrically aligned around the vertical dashed lines.
- The graph is likely illustrating a function over the interval from "a" to "b." 

The figure represents a common illustration in calculus or physics, possibly showing integrals or potential energy.

\#\# Dataset metadata

Electrostatics; Numerical Reasoning, Physical Model Grounding Reasoning

\paragraph{Recorded answer/context.}
D: D: 1

\paragraph{Figure/image path.}
/root/proposal\_for\_physic/science-mango/phyx\_mini\_run/phyx\_data/test\_image/862.png

\paragraph{Formalization target.}
create a compiling Lean file with sorry bodies at `PhyXMiniProblems/problem\_phyx\_mini\_0862.lean`. The Lean declarations must preserve the physical quantities, dimensions or dimensional roles, figure labels, governing-law hypotheses, and final relation expressed by this problem.
Use Mathlib/Physlib names found through LeanExplore where available. If a physics API is missing, introduce faithful local abstractions rather than scalar placeholder aliases.

\begin{theorem}[Physics formalization target]
\label{thm:physics:phyx_mini_0862:target}
\lean{PhyXMiniProblems.ProblemPhyXMini0862.problem_phyx_mini_0862}
\uses{def:physics:phyx-mini-0862:phyxminiproblems-problemphyxmini0862-chargeincoulombs, def:physics:phyx-mini-0862:phyxminiproblems-problemphyxmini0862-twopointchargexaxissetup, def:physics:phyx-mini-0862:phyxminiproblems-problemphyxmini0862-hassymmetricexteriorsamplegeometry, def:physics:phyx-mini-0862:phyxminiproblems-problemphyxmini0862-matchessuppliedelectricfieldgraph, def:physics:phyx-mini-0862:phyxminiproblems-problemphyxmini0862-traceheightcalibratesaxialfieldmagnitude}
The assigned autoformalize agent should translate this physics problem into Lean declarations in the covered file, with theorem and lemma proof bodies written as `by sorry`.
\end{theorem}
\begin{proof}
This is an autoformalization task, not a proof task. Produce statements that can later be proved by the physics prover without weakening the physical model.
\end{proof}

% --- Archon declaration topology begin ---
\section{Declaration topology}
\begin{definition}[electric Field Dimension]
  \label{def:physics:phyx-mini-0862:phyxminiproblems-problemphyxmini0862-electricfielddimension}
  \lean{PhyXMiniProblems.ProblemPhyXMini0862.electricFieldDimension}
  The physical dimension M L T⁻² C⁻¹ of an electric-field component
\end{definition}
\begin{proof}
  This is the declared model definition of electric Field Dimension.
\end{proof}
\begin{definition}[Signed Charge Quantity]
  \label{def:physics:phyx-mini-0862:phyxminiproblems-problemphyxmini0862-signedchargequantity}
  \lean{PhyXMiniProblems.ProblemPhyXMini0862.SignedChargeQuantity}
  A unit-independent signed electric charge
\end{definition}
\begin{proof}
  This is the declared model definition of Signed Charge Quantity.
\end{proof}
\begin{definition}[Axial Position Quantity]
  \label{def:physics:phyx-mini-0862:phyxminiproblems-problemphyxmini0862-axialpositionquantity}
  \lean{PhyXMiniProblems.ProblemPhyXMini0862.AxialPositionQuantity}
  A unit-independent signed position coordinate on the physical x-axis
\end{definition}
\begin{proof}
  This is the declared model definition of Axial Position Quantity.
\end{proof}
\begin{definition}[charge In Coulombs]
  \label{def:physics:phyx-mini-0862:phyxminiproblems-problemphyxmini0862-chargeincoulombs}
  \lean{PhyXMiniProblems.ProblemPhyXMini0862.chargeInCoulombs}
  \uses{def:physics:phyx-mini-0862:phyxminiproblems-problemphyxmini0862-signedchargequantity}
  Read a signed charge in coherent-SI coulombs
\end{definition}
\begin{proof}
  This is the declared model definition of charge In Coulombs.
\end{proof}
\begin{definition}[position In Meters]
  \label{def:physics:phyx-mini-0862:phyxminiproblems-problemphyxmini0862-positioninmeters}
  \lean{PhyXMiniProblems.ProblemPhyXMini0862.positionInMeters}
  \uses{def:physics:phyx-mini-0862:phyxminiproblems-problemphyxmini0862-axialpositionquantity}
  Read an axial position in coherent-SI metres
\end{definition}
\begin{proof}
  This is the declared model definition of position In Meters.
\end{proof}
\begin{definition}[x Axis Point In Meters]
  \label{def:physics:phyx-mini-0862:phyxminiproblems-problemphyxmini0862-xaxispointinmeters}
  \lean{PhyXMiniProblems.ProblemPhyXMini0862.xAxisPointInMeters}
  Embed a metre coordinate as a point on the first axis of physical 3-space
\end{definition}
\begin{proof}
  This is the declared model definition of x Axis Point In Meters.
\end{proof}
\begin{definition}[Charge Site]
  \label{def:physics:phyx-mini-0862:phyxminiproblems-problemphyxmini0862-chargesite}
  \lean{PhyXMiniProblems.ProblemPhyXMini0862.ChargeSite}
  The two sources named by the labels under the dashed lines
\end{definition}
\begin{proof}
  This is the declared model definition of Charge Site.
\end{proof}
\begin{definition}[expected Position Label]
  \label{def:physics:phyx-mini-0862:phyxminiproblems-problemphyxmini0862-expectedpositionlabel}
  \lean{PhyXMiniProblems.ProblemPhyXMini0862.expectedPositionLabel}
  \uses{def:physics:phyx-mini-0862:phyxminiproblems-problemphyxmini0862-chargesite}
  The text expected under each named source line
\end{definition}
\begin{proof}
  This is the declared model definition of expected Position Label.
\end{proof}
\begin{definition}[Graph Axis]
  \label{def:physics:phyx-mini-0862:phyxminiproblems-problemphyxmini0862-graphaxis}
  \lean{PhyXMiniProblems.ProblemPhyXMini0862.GraphAxis}
  The two semantic axes of the supplied graph
\end{definition}
\begin{proof}
  This is the declared model definition of Graph Axis.
\end{proof}
\begin{definition}[Graph Axis Role]
  \label{def:physics:phyx-mini-0862:phyxminiproblems-problemphyxmini0862-graphaxisrole}
  \lean{PhyXMiniProblems.ProblemPhyXMini0862.GraphAxisRole}
  The physical role assigned to a graph axis
\end{definition}
\begin{proof}
  This is the declared model definition of Graph Axis Role.
\end{proof}
\begin{definition}[Graph Branch]
  \label{def:physics:phyx-mini-0862:phyxminiproblems-problemphyxmini0862-graphbranch}
  \lean{PhyXMiniProblems.ProblemPhyXMini0862.GraphBranch}
  The three visible branches separated by the two source singularities
\end{definition}
\begin{proof}
  This is the declared model definition of Graph Branch.
\end{proof}
\begin{definition}[Electric Field Component Figure]
  \label{def:physics:phyx-mini-0862:phyxminiproblems-problemphyxmini0862-electricfieldcomponentfigure}
  \lean{PhyXMiniProblems.ProblemPhyXMini0862.ElectricFieldComponentFigure}
  \uses{def:physics:phyx-mini-0862:phyxminiproblems-problemphyxmini0862-chargesite, def:physics:phyx-mini-0862:phyxminiproblems-problemphyxmini0862-graphaxis, def:physics:phyx-mini-0862:phyxminiproblems-problemphyxmini0862-graphaxisrole, def:physics:phyx-mini-0862:phyxminiproblems-problemphyxmini0862-graphbranch}
  Literal presentation data and an idealized scalar-height readout for image 862.png. Heights are measured in the same numerical scale as newtons per coulomb; their connection to the physical field is stated separately
\end{definition}
\begin{proof}
  This is the declared model definition of Electric Field Component Figure.
\end{proof}
\begin{definition}[Two Point Charge XAxis Setup]
  \label{def:physics:phyx-mini-0862:phyxminiproblems-problemphyxmini0862-twopointchargexaxissetup}
  \lean{PhyXMiniProblems.ProblemPhyXMini0862.TwoPointChargeXAxisSetup}
  \uses{def:physics:phyx-mini-0862:phyxminiproblems-problemphyxmini0862-signedchargequantity, def:physics:phyx-mini-0862:phyxminiproblems-problemphyxmini0862-axialpositionquantity, def:physics:phyx-mini-0862:phyxminiproblems-problemphyxmini0862-chargesite, def:physics:phyx-mini-0862:phyxminiproblems-problemphyxmini0862-electricfieldcomponentfigure}
  The free-space medium, dimensionful source data, primary figure, and a pair of independent exterior sample positions used to read the pictured symmetry
\end{definition}
\begin{proof}
  This is the declared model definition of Two Point Charge XAxis Setup.
\end{proof}
\begin{definition}[source Point Charge Potential]
  \label{def:physics:phyx-mini-0862:phyxminiproblems-problemphyxmini0862-sourcepointchargepotential}
  \lean{PhyXMiniProblems.ProblemPhyXMini0862.sourcePointChargePotential}
  \uses{def:physics:phyx-mini-0862:phyxminiproblems-problemphyxmini0862-chargeincoulombs, def:physics:phyx-mini-0862:phyxminiproblems-problemphyxmini0862-positioninmeters, def:physics:phyx-mini-0862:phyxminiproblems-problemphyxmini0862-xaxispointinmeters, def:physics:phyx-mini-0862:phyxminiproblems-problemphyxmini0862-chargesite, def:physics:phyx-mini-0862:phyxminiproblems-problemphyxmini0862-twopointchargexaxissetup}
  The stationary 3D point-particle potential attached to a named source
\end{definition}
\begin{proof}
  This is the declared model definition of source Point Charge Potential.
\end{proof}
\begin{definition}[net Point Charge Potential]
  \label{def:physics:phyx-mini-0862:phyxminiproblems-problemphyxmini0862-netpointchargepotential}
  \lean{PhyXMiniProblems.ProblemPhyXMini0862.netPointChargePotential}
  \uses{def:physics:phyx-mini-0862:phyxminiproblems-problemphyxmini0862-chargesite, def:physics:phyx-mini-0862:phyxminiproblems-problemphyxmini0862-twopointchargexaxissetup, def:physics:phyx-mini-0862:phyxminiproblems-problemphyxmini0862-sourcepointchargepotential}
  The sum of the two stationary point-particle potentials
\end{definition}
\begin{proof}
  This is the declared model definition of net Point Charge Potential.
\end{proof}
\begin{definition}[physlib Source Point Charge Electric Field Law]
  \label{def:physics:phyx-mini-0862:phyxminiproblems-problemphyxmini0862-physlibsourcepointchargeelectricfieldlaw}
  \lean{PhyXMiniProblems.ProblemPhyXMini0862.physlibSourcePointChargeElectricFieldLaw}
  \uses{def:physics:phyx-mini-0862:phyxminiproblems-problemphyxmini0862-chargeincoulombs, def:physics:phyx-mini-0862:phyxminiproblems-problemphyxmini0862-positioninmeters, def:physics:phyx-mini-0862:phyxminiproblems-problemphyxmini0862-xaxispointinmeters, def:physics:phyx-mini-0862:phyxminiproblems-problemphyxmini0862-chargesite, def:physics:phyx-mini-0862:phyxminiproblems-problemphyxmini0862-twopointchargexaxissetup}
  An exact specialization of Physlib's theorem giving the distributional electric field of each source. This declaration records the library law used to ground the ordinary axial readout below; it introduces no new axiom
\end{definition}
\begin{proof}
  This is the declared model definition of physlib Source Point Charge Electric Field Law.
\end{proof}
\begin{definition}[source Point Charge Electric Field Ordinary Profile]
  \label{def:physics:phyx-mini-0862:phyxminiproblems-problemphyxmini0862-sourcepointchargeelectricfieldordinaryprofile}
  \lean{PhyXMiniProblems.ProblemPhyXMini0862.sourcePointChargeElectricFieldOrdinaryProfile}
  \uses{def:physics:phyx-mini-0862:phyxminiproblems-problemphyxmini0862-chargeincoulombs, def:physics:phyx-mini-0862:phyxminiproblems-problemphyxmini0862-positioninmeters, def:physics:phyx-mini-0862:phyxminiproblems-problemphyxmini0862-xaxispointinmeters, def:physics:phyx-mini-0862:phyxminiproblems-problemphyxmini0862-chargesite, def:physics:phyx-mini-0862:phyxminiproblems-problemphyxmini0862-twopointchargexaxissetup}
  The ordinary vector profile occurring inside Physlib's distributional point-charge field theorem, with coherent-SI source data
\end{definition}
\begin{proof}
  This is the declared model definition of source Point Charge Electric Field Ordinary Profile.
\end{proof}
\begin{definition}[coulomb Constant In Newton Meters Squared Per Coulomb Squared]
  \label{def:physics:phyx-mini-0862:phyxminiproblems-problemphyxmini0862-coulombconstantinnewtonmeterssquaredpercoulombsquared}
  \lean{PhyXMiniProblems.ProblemPhyXMini0862.coulombConstantInNewtonMetersSquaredPerCoulombSquared}
  \uses{def:physics:phyx-mini-0862:phyxminiproblems-problemphyxmini0862-twopointchargexaxissetup}
  The free-space Coulomb constant in coherent SI, N m² / C²
\end{definition}
\begin{proof}
  This is the declared model definition of coulomb Constant In Newton Meters Squared Per Coulomb Squared.
\end{proof}
\begin{definition}[displacement From Source In Meters]
  \label{def:physics:phyx-mini-0862:phyxminiproblems-problemphyxmini0862-displacementfromsourceinmeters}
  \lean{PhyXMiniProblems.ProblemPhyXMini0862.displacementFromSourceInMeters}
  \uses{def:physics:phyx-mini-0862:phyxminiproblems-problemphyxmini0862-axialpositionquantity, def:physics:phyx-mini-0862:phyxminiproblems-problemphyxmini0862-positioninmeters, def:physics:phyx-mini-0862:phyxminiproblems-problemphyxmini0862-chargesite, def:physics:phyx-mini-0862:phyxminiproblems-problemphyxmini0862-twopointchargexaxissetup}
  Metre displacement from a source to an axial observation point
\end{definition}
\begin{proof}
  This is the declared model definition of displacement From Source In Meters.
\end{proof}
\begin{definition}[source Axial Electric Field In Newtons Per Coulomb]
  \label{def:physics:phyx-mini-0862:phyxminiproblems-problemphyxmini0862-sourceaxialelectricfieldinnewtonspercoulomb}
  \lean{PhyXMiniProblems.ProblemPhyXMini0862.sourceAxialElectricFieldInNewtonsPerCoulomb}
  \uses{def:physics:phyx-mini-0862:phyxminiproblems-problemphyxmini0862-axialpositionquantity, def:physics:phyx-mini-0862:phyxminiproblems-problemphyxmini0862-positioninmeters, def:physics:phyx-mini-0862:phyxminiproblems-problemphyxmini0862-xaxispointinmeters, def:physics:phyx-mini-0862:phyxminiproblems-problemphyxmini0862-chargesite, def:physics:phyx-mini-0862:phyxminiproblems-problemphyxmini0862-twopointchargexaxissetup, def:physics:phyx-mini-0862:phyxminiproblems-problemphyxmini0862-sourcepointchargeelectricfieldordinaryprofile}
  The signed coherent-SI x-component of one source's field. It is the x-axis restriction of Physlib's three-dimensional point-particle profile q r / ‖r‖³
\end{definition}
\begin{proof}
  This is the declared model definition of source Axial Electric Field In Newtons Per Coulomb.
\end{proof}
\begin{lemma}[source axial profile eq coulomb formula]
  \label{lem:physics:phyx-mini-0862:phyxminiproblems-problemphyxmini0862-source-axial-profile-eq-coulomb-formula}
  \lean{PhyXMiniProblems.ProblemPhyXMini0862.source_axial_profile_eq_coulomb_formula}
  \uses{def:physics:phyx-mini-0862:phyxminiproblems-problemphyxmini0862-axialpositionquantity, def:physics:phyx-mini-0862:phyxminiproblems-problemphyxmini0862-chargeincoulombs, def:physics:phyx-mini-0862:phyxminiproblems-problemphyxmini0862-chargesite, def:physics:phyx-mini-0862:phyxminiproblems-problemphyxmini0862-twopointchargexaxissetup, def:physics:phyx-mini-0862:phyxminiproblems-problemphyxmini0862-coulombconstantinnewtonmeterssquaredpercoulombsquared, def:physics:phyx-mini-0862:phyxminiproblems-problemphyxmini0862-displacementfromsourceinmeters, def:physics:phyx-mini-0862:phyxminiproblems-problemphyxmini0862-sourceaxialelectricfieldinnewtonspercoulomb}
  On the x-axis, the Physlib ordinary vector profile reduces to the familiar signed scalar form k q (x-x₀) / |x-x₀|³
\end{lemma}
\begin{proof}
  The conclusion follows from the governing relations and figure readouts named in the statement of source axial profile eq coulomb formula.
\end{proof}
\begin{definition}[net Axial Electric Field In Newtons Per Coulomb]
  \label{def:physics:phyx-mini-0862:phyxminiproblems-problemphyxmini0862-netaxialelectricfieldinnewtonspercoulomb}
  \lean{PhyXMiniProblems.ProblemPhyXMini0862.netAxialElectricFieldInNewtonsPerCoulomb}
  \uses{def:physics:phyx-mini-0862:phyxminiproblems-problemphyxmini0862-axialpositionquantity, def:physics:phyx-mini-0862:phyxminiproblems-problemphyxmini0862-chargesite, def:physics:phyx-mini-0862:phyxminiproblems-problemphyxmini0862-twopointchargexaxissetup, def:physics:phyx-mini-0862:phyxminiproblems-problemphyxmini0862-sourceaxialelectricfieldinnewtonspercoulomb}
  Linear superposition of the two signed axial point-charge fields
\end{definition}
\begin{proof}
  This is the declared model definition of net Axial Electric Field In Newtons Per Coulomb.
\end{proof}
\begin{definition}[Has Symmetric Exterior Sample Geometry]
  \label{def:physics:phyx-mini-0862:phyxminiproblems-problemphyxmini0862-hassymmetricexteriorsamplegeometry}
  \lean{PhyXMiniProblems.ProblemPhyXMini0862.HasSymmetricExteriorSampleGeometry}
  \uses{def:physics:phyx-mini-0862:phyxminiproblems-problemphyxmini0862-chargeincoulombs, def:physics:phyx-mini-0862:phyxminiproblems-problemphyxmini0862-positioninmeters, def:physics:phyx-mini-0862:phyxminiproblems-problemphyxmini0862-twopointchargexaxissetup}
  The charges are distinct and nonzero, and the two chosen samples are exterior points reflected through the midpoint of a and b. These hypotheses make the requested ratio well-defined without assuming its value
\end{definition}
\begin{proof}
  This is the declared model definition of Has Symmetric Exterior Sample Geometry.
\end{proof}
\begin{definition}[Matches Supplied Electric Field Graph]
  \label{def:physics:phyx-mini-0862:phyxminiproblems-problemphyxmini0862-matchessuppliedelectricfieldgraph}
  \lean{PhyXMiniProblems.ProblemPhyXMini0862.MatchesSuppliedElectricFieldGraph}
  \uses{def:physics:phyx-mini-0862:phyxminiproblems-problemphyxmini0862-positioninmeters, def:physics:phyx-mini-0862:phyxminiproblems-problemphyxmini0862-expectedpositionlabel, def:physics:phyx-mini-0862:phyxminiproblems-problemphyxmini0862-twopointchargexaxissetup}
  Literal raster features plus the idealized equal-height readout at the chosen mirror pair. This predicate contains no charge or charge-ratio equality
\end{definition}
\begin{proof}
  This is the declared model definition of Matches Supplied Electric Field Graph.
\end{proof}
\begin{definition}[Trace Height Calibrates Axial Field Magnitude]
  \label{def:physics:phyx-mini-0862:phyxminiproblems-problemphyxmini0862-traceheightcalibratesaxialfieldmagnitude}
  \lean{PhyXMiniProblems.ProblemPhyXMini0862.TraceHeightCalibratesAxialFieldMagnitude}
  \uses{def:physics:phyx-mini-0862:phyxminiproblems-problemphyxmini0862-positioninmeters, def:physics:phyx-mini-0862:phyxminiproblems-problemphyxmini0862-twopointchargexaxissetup, def:physics:phyx-mini-0862:phyxminiproblems-problemphyxmini0862-netaxialelectricfieldinnewtonspercoulomb}
  The prose identifies the plotted profile with E\_x; because the raster has only nonnegative heights, its ordinate is calibrated to the magnitude of the signed net component. This is a figure-to-physics calibration, not an answer premise
\end{definition}
\begin{proof}
  This is the declared model definition of Trace Height Calibrates Axial Field Magnitude.
\end{proof}
\begin{lemma}[symmetric sample field magnitudes equal]
  \label{lem:physics:phyx-mini-0862:phyxminiproblems-problemphyxmini0862-symmetric-sample-field-magnitudes-equal}
  \lean{PhyXMiniProblems.ProblemPhyXMini0862.symmetric_sample_field_magnitudes_equal}
  \uses{def:physics:phyx-mini-0862:phyxminiproblems-problemphyxmini0862-twopointchargexaxissetup, def:physics:phyx-mini-0862:phyxminiproblems-problemphyxmini0862-netaxialelectricfieldinnewtonspercoulomb, def:physics:phyx-mini-0862:phyxminiproblems-problemphyxmini0862-matchessuppliedelectricfieldgraph, def:physics:phyx-mini-0862:phyxminiproblems-problemphyxmini0862-traceheightcalibratesaxialfieldmagnitude}
  The mirror sample readouts give equal magnitudes of the net axial field
\end{lemma}
\begin{proof}
  The conclusion follows from the governing relations and figure readouts named in the statement of symmetric sample field magnitudes equal.
\end{proof}
\begin{lemma}[charge readout magnitudes equal]
  \label{lem:physics:phyx-mini-0862:phyxminiproblems-problemphyxmini0862-charge-readout-magnitudes-equal}
  \lean{PhyXMiniProblems.ProblemPhyXMini0862.charge_readout_magnitudes_equal}
  \uses{def:physics:phyx-mini-0862:phyxminiproblems-problemphyxmini0862-chargeincoulombs, def:physics:phyx-mini-0862:phyxminiproblems-problemphyxmini0862-twopointchargexaxissetup, def:physics:phyx-mini-0862:phyxminiproblems-problemphyxmini0862-hassymmetricexteriorsamplegeometry, def:physics:phyx-mini-0862:phyxminiproblems-problemphyxmini0862-matchessuppliedelectricfieldgraph, def:physics:phyx-mini-0862:phyxminiproblems-problemphyxmini0862-traceheightcalibratesaxialfieldmagnitude}
  For two nonzero point charges at mirror exterior samples, the equal field magnitudes force equality of the source-charge magnitudes
\end{lemma}
\begin{proof}
  The conclusion follows from the governing relations and figure readouts named in the statement of charge readout magnitudes equal.
\end{proof}
\begin{definition}[Answer Choice]
  \label{def:physics:phyx-mini-0862:phyxminiproblems-problemphyxmini0862-answerchoice}
  \lean{PhyXMiniProblems.ProblemPhyXMini0862.AnswerChoice}
  Labels attached to the four ratio choices in the source
\end{definition}
\begin{proof}
  This is the declared model definition of Answer Choice.
\end{proof}
\begin{definition}[ratio]
  \label{def:physics:phyx-mini-0862:phyxminiproblems-problemphyxmini0862-answerchoice-ratio}
  \lean{PhyXMiniProblems.ProblemPhyXMini0862.AnswerChoice.ratio}
  \uses{def:physics:phyx-mini-0862:phyxminiproblems-problemphyxmini0862-answerchoice}
  The dimensionless ratio printed beside each choice
\end{definition}
\begin{proof}
  This is the declared model definition of ratio.
\end{proof}
\begin{definition}[recorded Dataset Answer]
  \label{def:physics:phyx-mini-0862:phyxminiproblems-problemphyxmini0862-recordeddatasetanswer}
  \lean{PhyXMiniProblems.ProblemPhyXMini0862.recordedDatasetAnswer}
  \uses{def:physics:phyx-mini-0862:phyxminiproblems-problemphyxmini0862-answerchoice}
  The answer label recorded in the supplied dataset metadata
\end{definition}
\begin{proof}
  This is the declared model definition of recorded Dataset Answer.
\end{proof}
% --- Archon declaration topology end ---
% --- Archon physics formalization source end ---
